\section{Analysis and Discussion}
\label{sec:analysis}

\subsection{Probe Instability and Ensemble Resolution}
\label{sec:analysis:instability}

Our most important methodological finding is that single-seed,
single-split probe training on small datasets (${\sim}$2{,}600
samples, 15\% validation $\approx$ 400 samples) produces highly
unstable $R^2$ scores. Under naive training, two probes failed entirely
(LTV/Loan: $R^2 = -0.03$; Debt Yield/NOI: $R^2 = 0.46$), while
others varied by up to 0.75 across random seeds.

\begin{table}[h]
\centering
\begin{tabular}{@{}lccc@{}}
\toprule
\textbf{Expert} & \textbf{Single-seed $R^2$} & \textbf{5-fold CV $R^2$} & \textbf{$<$20\% Acc.} \\
\midrule
DSCR       & 0.97  & $0.885 \pm 0.183$ & 60\% \\
Cap Rate   & 0.97  & $0.975 \pm 0.002$ & 65\% \\
LTV        & $-$0.03 & $0.982 \pm 0.001$ & 87\% \\
Debt Yield & 0.46  & $0.977 \pm 0.002$ & 65\% \\
\bottomrule
\end{tabular}
\caption{Single-seed training vs.\ 5-fold CV ensemble for the
worst-case input probe per expert. LTV and Debt Yield, which
completely failed under single-seed training, are fully recovered
by the ensemble. DSCR/NOI remains the most variable probe
($0.885 \pm 0.183$) but the 5-probe ensemble median suppresses
the outlier fold at inference.}
\label{tab:bottleneck}
\end{table}

The 5-fold CV ensemble resolves this: by training 5 probes on
different data splits (each with 3 random seeds, keeping the best),
the ensemble median suppresses individual probe failures. This
demonstrates that the extraction failures were \emph{not}
fundamental limits of the hidden state representation but rather
artifacts of overfitting to a particular train/val split.

\subsection{Error Compounding in Chained Computation}
\label{sec:analysis:compounding}

Deterministic computation amplifies extraction errors because there is
no opportunity for neural compensation. For DSCR = NOI / ADS, if both
inputs have ${\sim}$10\% error, the quotient can have ${\sim}$20\%
error. This explains why DSCR has the highest median error (15.5\%)
despite both input probes achieving $R^2 > 0.88$.

Direct output probes, which predict the final metric value directly
from hidden states bypassing the deterministic computation, achieve
lower median errors (12.3\% vs.\ 15.5\% for DSCR) despite
substantially lower $R^2$ ($0.46$--$0.56$ for DSCR, Cap Rate, and
Debt Yield; LTV direct output $R^2$ is negative, yet still matches
the ensemble pipeline at 86.5\% accuracy). This suggests that
for division-based metrics, the ``extract-then-compute'' pipeline
is suboptimal compared to direct prediction, and that the hidden
state encodes the \emph{answer} (not just the inputs) at a moderate
but useful level of fidelity.

\subsection{Log-Scale Training and Error Amplification}
\label{sec:analysis:logscale}

We train extraction probes in log-scale ($\text{sign}(y) \cdot
\log(1 + |y|)$) to handle the 1{,}000$\times$ dynamic range of
financial values. While this enables learning across the full range,
it introduces multiplicative error when converting back to raw values.
A residual standard deviation of 0.24 in log-space translates to
approximately $e^{0.24} - 1 \approx 27\%$ multiplicative error in
raw values. This partially explains median errors of 6--16\%: the
extraction is strong in log-space ($R^2 > 0.97$) but the
nonlinear back-transformation amplifies small errors.

\subsection{Layer Selection}
\label{sec:analysis:layer}

We evaluate extraction at layers 12, 16, and 20 (of 24 total). Layer~12
consistently outperforms deeper layers for input value extraction. This
is consistent with prior work showing that factual/numerical information
is most accessible in middle layers
\citep{geva2023dissecting,meng2022locating}, while later layers shift
toward next-token prediction representations.

\subsection{Comparison to Tool-Calling and Generation Approaches}
\label{sec:analysis:toolcalling}

Our approach trades per-prediction accuracy for \emph{coverage} and
\emph{speed}:

\begin{itemize}
\item \textbf{Generation baselines} achieve high per-prediction accuracy
when they produce parseable output (e.g., 96--100\% for structured
Cap Rate, Debt Yield) but fail to parse a significant fraction of
scenarios ($N$ ranges from 0 to 187 out of 200). The structured prompt
improves coverage but degrades accuracy on hard metrics (DSCR drops
from 77\% to 19\%).
\item \textbf{Tool calling} achieves perfect computation accuracy
\emph{when invoked correctly}, but requires multiple inference passes
(2--5$\times$ total latency) and an external execution environment.
\item \textbf{Hybrid MoE} guarantees a prediction for every scenario
($N{=}200$) at 173$\times$ the speed, with 60--87\% within 20\%
accuracy. Computation is \emph{exactly} correct; accuracy is bounded
solely by extraction quality.
\end{itemize}

This positions Hybrid MoE as ideal for high-throughput, latency-sensitive
applications (batch underwriting, real-time screening) where guaranteed
coverage at 173$\times$ the speed is preferable to high per-prediction
accuracy with incomplete coverage and multi-second latency.

\subsection{Limitations}
\label{sec:analysis:limitations}

\begin{enumerate}
\item \textbf{Extraction is the bottleneck.} Accuracy depends entirely
on extraction quality. Despite $R^2 \geq 0.88$ for all probes,
median errors of 6--16\% persist due to log-scale
back-transformation and error compounding through division.
Attention-based extraction over specific token positions is the
primary avenue for improvement.

\item \textbf{Probe seed sensitivity.} On small training sets
(${\sim}$2{,}600 samples), individual probes are sensitive to random
initialization and train/val splits. Our 5-fold CV ensemble
mitigates this but adds training complexity. Larger training sets
would reduce this sensitivity.

\item \textbf{Baselines are under-optimized.} Our generation baselines
use a base model without instruction tuning. An instruction-tuned
variant would likely produce more consistently formatted outputs,
increasing baseline $N$ and potentially accuracy. The structured
prompt partially addresses this.

\item \textbf{Synthetic evaluation.} Our CRE benchmark is synthetically
generated. Real-world underwriting documents contain ambiguity,
missing data, and domain-specific conventions not captured here.

\item \textbf{Mean-pooling limitation.} We extract from mean-pooled
hidden states, which loses positional information about where
specific numbers appear in the input.

\item \textbf{Limited expert set.} We evaluate four CRE metrics. The
architecture supports arbitrary Python functions, but generalization
to more complex formulas was not validated.

\item \textbf{Standalone probes, not integrated forward pass.} Our
experiments train extraction MLPs on cached hidden states and use a
separate binary classifier for routing, rather than running the
extended router end-to-end within the MoE layer. This validates
that the required information is extractable and the components
work, but does not test the fully integrated architecture
(including projection back into the residual stream and router
competition between neural and deterministic experts). End-to-end
integration is future work.
\end{enumerate}
